% !TeX root = main.tex
\section{Introducción}
\label{sec:introduccion}

La segmentación de imágenes asigna etiquetas a los píxeles para delimitar
regiones de interés~\cite{szeliski2011}. En una imagen de resonancia magnética, la delimitación
de una lesión requiere distinguirla de los tejidos circundantes a partir
de su apariencia y de su organización espacial. Cuando distintas regiones
presentan intensidades semejantes, una clasificación basada únicamente en
la intensidad puede producir ambigüedades. Incorporar relaciones entre
píxeles vecinos permite favorecer regiones coherentes, mientras que una
selección interactiva restringe la búsqueda a la zona de interés~\cite{boykov2001}.

GrabCut combina modelos de apariencia de fondo y primer plano con una
penalización espacial dependiente del contraste. El procedimiento alterna
la estimación de estos modelos con la actualización de etiquetas mediante
cortes de grafos. Una inicialización rectangular permite indicar la región
que contiene el objeto sin trazar de antemano su borde completo~\cite{rother2004}.
La segmentación obtenida depende tanto de la información de la imagen como
de la restricción introducida por ese rectángulo.

En este trabajo se estudia una implementación de segmentación interactiva
basada en GrabCut aplicada a cortes de resonancia magnética de schwannoma
vestibular. Se utiliza un conjunto de trabajo de 30 imágenes PNG asociado
a la colección \textit{Vestibular-Schwannoma-SEG}~\cite{shapeydatos}, cuya procedencia y
preparación se describen en la sección~\ref{sec:datos}. Las imágenes se
procesan individualmente en escala de grises. Para cada corte, la región
candidata se delimita mediante un rectángulo y se construyen modelos de
mezclas gaussianas para ambas clases. La coherencia espacial se incorpora
mediante una vecindad de ocho píxeles y las etiquetas se actualizan por
corte mínimo. La anotación roja disponible se transforma en una máscara
binaria destinada a la evaluación.

El objetivo es describir la relación entre la inicialización interactiva,
los modelos de intensidad y la optimización espacial, y examinar la
segmentación resultante mediante comparaciones con la referencia. Se
registra la evolución de las etiquetas y de la energía durante cinco
iteraciones. La evaluación considera el error cuadrático medio sobre la
imagen completa y sobre el rectángulo, los falsos positivos y negativos y
su distribución espacial. Se incluye una predicción de todo fondo para
interpretar el efecto del predominio de esta clase en el error global.

El procedimiento se ilustra con el caso \texttt{VS-SEG-001}, para el que
se presentan la evolución de la segmentación y la comparación con su
anotación de referencia.
